%% source: 2023-sp-midterm_01
\begin{prob}[(2 points)]
    Solve the below recurrence, stating the solution in asymptotic notation. Show your work.
    \[
        T(n) =
        \begin{cases}
            T(n/2) + \Theta(n) & n > 1 \\
            \Theta(1) & n = 1
        \end{cases}
    \]

    \begin{responsebox}{3.75in}
        $\Theta(n)$.

        Unrolling several times, we find:

        \begin{align*}
            T(n)
            &=
                T(n/2) + n
                \\
            &=
                T(n/4) + \frac{n}{2} + n
                \\
            &=
                T(n/8) + \frac{n}{4} + \frac{n}{2} + n
                \\
        \end{align*}

        On the $k$th unroll, we'll get:
        \[
            T(n) = T(n/2^k) + \left[
                n + \frac{n}{2} + \frac{n}{4} + \ldots + \frac{n}{2^{k-1}}
            \right]
        \]

        It will take $\log_2 n$ unrolls to each the base case. Substituting this for
        $k$, we get:
        \begin{align*}
            T(n)
            &= T(n/2^{\log_2 n}) + \left[
                n + \frac{n}{2} + \frac{n}{4} + \ldots + \frac{n}{2^{(\log_2 n)-1}}
            \right]
            \\
            &= T(1) + n
            \underbrace{\left[
                1 + \frac{1}{2} + \frac{1}{4} + \ldots + \frac{1}{2^{(\log_2 n)-1}}
            \right]}_{\leq 2}
        \intertext{Here we've recognized the summation in square braces as a geometric
            sum of powers of $1/2$, which converges to 2 when there are infinitely-many
            terms. We've seen this sum in solving a couple of homework problems.}
            & = T(1) + c n\\
            & = \Theta(1) + c n\\
            &= \Theta(n)
        \end{align*}

    \end{responsebox}
\end{prob}
